% \definecolor{bordergray}{RGB}{180, 180, 180}
% \definecolor{headerblue}{RGB}{220, 230, 240}
% \definecolor{paleblue}{RGB}{240, 248, 255}
% \definecolor{codekey}{RGB}{0, 0, 255}
% \definecolor{codestr}{RGB}{163, 21, 21}
% \definecolor{flowred}{RGB}{220, 50, 47}
% \definecolor{flowblue}{RGB}{38, 139, 210}
% \definecolor{llvmgreen}{RGB}{0, 100, 0}

% \tikzset{
%     % 基础面板
%     panel_base/.style={
%         draw=bordergray,
%         thick,
%         fill=white,
%         inner sep=3pt,
%         align=left,
%         anchor=north west,
%         outer sep=0pt
%     },
%     panel_slim/.style={ panel_base, text width=2.8cm, inner sep=2pt, outer sep=0pt},
%     panel_narrow/.style={ panel_base, text width=4.0cm, inner sep=2pt, outer sep=0pt},
%     panel_wide/.style={ panel_base, text width=5.2cm, inner sep=2pt, outer sep=0pt},
%     % 虚线背景面板 (去掉了 fill=paleblue，实现无底纹)
%     panel_dashed/.style={
%         draw=bordergray,
%         thick,
%         dashed,
%         % fill=paleblue, % 已注释掉背景色
%         inner sep=2pt,
%         align=left,
%         anchor=north west,
%         outer sep=0pt
%     },
%     % 标题栏样式
%     header/.style={
%         fill=headerblue,
%         draw=bordergray,
%         thick,
%         align=center,
%         font=\bfseries\scriptsize,
%         yshift=-0.5pt,
%         outer sep=0pt,
%         inner sep=2pt,
%         anchor=south west
%     },
%     % 连线样式
%     flowline/.style={->, >=stealth, thick, rounded corners=2pt},
%     typearrow/.style={->, >=stealth, dashed, thick, color=gray},
%     % 变量高亮框
%     boxstyle/.style={draw, thick, rounded corners=1pt}
% }

% \newcommand{\tightnode}[3][]{%
%     \tikzmarknode[%
%         inner xsep=1.5pt,
%         inner ysep=0.5pt,
%         text height=6.0pt,
%         text depth=1.5pt,
%         #1  % <-- 允许插入额外的样式覆盖默认设置
%     ]{#2}{%
%         \smash{#3}
%     }%
% }

% \lstset{
%     language=C++,
%     basicstyle=\ttfamily\fontsize{5pt}{6pt}\selectfont,
%     breaklines=true,
%     escapechar=@, 
%     numbers=left,
%     numberstyle=\tiny\color{gray!50},
%     numbersep=2pt,
%     keywordstyle=\color{codekey}\bfseries,
%     commentstyle=\color{gray},
%     stringstyle=\color{codestr},
%     showstringspaces=false,
%     frame=none,
%     xleftmargin=0.8em,
%     lineskip=-0.5pt,      
%     aboveskip=2pt,        
%     belowskip=0pt
% }

% \begin{figure}[htbp]
%     \centering
%     \begin{tikzpicture}[remember picture]
        
%         % ==========================
%         % 1. PyTorch Frontend (Python)
%         % ==========================
%         \node[panel_slim, text width=2.8cm] (p1) {
%         % --- Part A: ATen IR (Compiler View) ---
%         \textbf{\scriptsize\sffamily\color{black} [Compile Time / ATen IR]} \\
%         \begin{lstlisting}[language=llvm]
% %B: SymInt
% %S: SymInt
% %lens: Array(SymInt)
% %q: Tensor[%B,@\tightnode{initial_const}{40}@,%S,64]
% %out = aten::fa(%q, 
%   %k, %v, %lens, ...)
%         \end{lstlisting}
        
%         \vspace{2pt}
%         \color{gray}\hrule height 0.5pt \normalcolor
%         \vspace{2pt}
        
%         % --- Part B: Python API (User View) ---
%         \textbf{\scriptsize\sffamily\color{black} [Runtime / Op Invocation]} \\
%         \begin{lstlisting}[language=Python]
% while greedy_search:
%   ...
%   # lens grows
%   torch.fa(q_cur, k_cache, v_cache, @\tightnode{py_seq}{lens}@, ...)
%         \end{lstlisting}
%     };
%         \node[header, text width=2.8cm] at (p1.north west) {1. PyTorch (Python)};
        
%         % ==========================
%         % 2. PyBind Schema (C++)
%         % ==========================
%         \node[panel_slim, text width=2.8cm, right=0.3cm of p1.north east, anchor=north west] (p2) {
%      \begin{lstlisting}[language=C++]
% struct Tensor { 
%   void* data_;
%   int* sizes; 
%   int ndim; 
% };
% m.def("fa", [](
%   Tensor @\tightnode{pb_q}{q}@, ..., 
%   int[] @\tightnode{pb_seq}{s}@, ...) {
%   return host_fa(q, k, v, s, ...);
% });
%     \end{lstlisting}
%         };
%         \node[header, text width=2.8cm] at (p2.north west) {2. PyBind (C++)};

%         % ==========================
%         % 3. Host Runtime (C++)
%         % ==========================
%         \node[panel_narrow, text width=4.0cm, right=0.3cm of p2.north east, anchor=north west] (p3) {
%         \begin{lstlisting}
% Tensor host_fa(Tensor q, ...){
%  int h = @\tightnode{host_h}{q}@.size(1); // 40
%  int stages = (@\tightnode{host_stages}{h}@>32)? @\tightnode{llvm_s_val}{2}@ : 1;
%  EXEC(kernel, ..., stages, @\tightnode{host_lens}{sum(s)}@);
% }
%         \end{lstlisting}
%         };
%         \node[header, text width=4.0cm] at (p3.north west) {3. Host (C++)};

%         % ==========================
%         % 4. Device Kernel
%         % ==========================
%         \node[panel_wide, text width=4.0cm, right=0.3cm of p3.north east, anchor=north west] (p4) {
%             \begin{lstlisting}
% void kernel(..., int stages, int len_sum) {
%   if (tid >= @\tightnode{kern_lens}{len\_sum}@) return;
%   for(int i=0; i < @\tightnode{kern_stages}{stages}@; i++)
%     ...
% }
%             \end{lstlisting}
%         };
%         \node[header, text width=4.0cm] at (p4.north west) {4a. Generic Kernel (C++)};

%         \node[panel_wide, text width=3.6cm, below=0.8cm of p4.south east, anchor=north east] (p_spec) {
%         \begin{lstlisting}
% void kernel(..., int len_sum) {
%   if (tid >= @\tightnode{kern_lens_use}{len\_sum}@) return; 
%   @\tightnode[inner sep=3pt]{kern_stage_spec}{%
%     \begin{tabular}{@{}l@{}}
%       ... // Stage 1\\
%       ... // Stage 2
%     \end{tabular}%
%   }@
% }
%         \end{lstlisting}
%         };
%         \node[header, text width=3.6cm] (p_spec_header) at (p_spec.north west) {4b. Specialized Kernel (C++)};

%         % ==========================
%         % 3.5 LLVM IR (JIT Analysis)
%         % ==========================
%         \node[panel_dashed, text width=4.4cm, below=0.8cm of p3.south west, anchor=north west] (p_llvm) {
%         \begin{lstlisting}[language=llvm]
% ...
% sizes = getelementptr %q ...
% h_ptr = getelementptr %sizes, i32 4
% h_val = load i32, %h_ptr  ; Load 40
% cmp = icmp sgt i32 %h_val, 32
% stages = select i1 %cmp, i32 2, i32 1
%         \end{lstlisting}
%         };
%         \node[header, text width=4.4cm] (p_llvm_header) at (p_llvm.north west) {3. Propagation on LLVM IR of Host};

% %         % ==========================
% %         % Type Mapping Layer - MODIFIED
% %         % ==========================
% %         \node[panel_dashed, text width=6cm, anchor=north west] (p_typemap) at (p1.south west |- p_llvm.north west) {
% %             \vspace{2pt} % Leave space for header
% %             \begin{minipage}{0.35\textwidth}
% %             \begin{lstlisting}[basicstyle=\ttfamily\tiny, numbers=none]
% % // C++ Layout            
% % struct Tensor {
% %  void* data_;
% %  int* sizes;
% %  int ndim;
% % };
% %             \end{lstlisting}
% %             \end{minipage}%
% %             \hfill
% %             \begin{minipage}{0.63\textwidth}
% %             \begin{lstlisting}[basicstyle=\ttfamily\tiny, numbers=none]
% % // PyBind Registration
% % py::class_<Tensor>(m, "Tensor")
% %   .def("size", [](Tensor& t, int id){ 
% %     return t.sizes[id]; 
% %  });
% %             \end{lstlisting}
% %             \end{minipage}
% %         };
% %         % Added Header for Type Mapping Layer
% %         \node[header, text width=6cm] (p_typemap_header) at (p_typemap.north west) {Type Mapping Layer (Python $\leftrightarrow$ C++)};

%     \end{tikzpicture}

%     % --- 绘制连线 (Overlay) ---
%     \begin{tikzpicture}[remember picture, overlay]
        
%         % --- RED FLOW: Metadata/Constant Flow ---
%         % 1. Python (40) SE -> PyBind (q) NW
%         \draw[flowred, boxstyle] (initial_const.north west) rectangle (initial_const.south east);
%         \draw[flowred, boxstyle] (pb_q.north west) rectangle (pb_q.south east);
%         \draw[flowline, flowred] (initial_const.north east) -- ++(1,0) |- (pb_q.north west);
        
%         % 2. PyBind (q) SE -> Host (q) NW
%         \draw[flowred, boxstyle] (host_h.north west) rectangle (host_h.south east);
%         \draw[flowline, flowred] (pb_q.north east) -- ++(1,0) |- (host_h.north west);
        
%         % 3. Host (q) SW -> Host (S) NE
%         \draw[flowred, boxstyle] (host_stages.north west) rectangle (host_stages.south east);
%         \draw[flowline, flowred] (host_h.south west) -- (host_stages.north west);

%         % 4. Host (S) -> LLVM / Device (IR Generation)
%         % \draw[flowline, flowred, dashed, opacity=0.5] (host_stages.south) -- (p_llvm_header.north);

%         % 5. LLVM -> Device Specialization
%         \draw[flowred, boxstyle] (llvm_s_val.north west) rectangle (llvm_s_val.south east);
%         \draw[flowred, boxstyle] (kern_stages.north west) rectangle (kern_stages.south east);
%         \draw[flowline, flowred] (llvm_s_val.south east) -- ++(1,0) |- (kern_stages.south west);

%         % --- BLUE FLOW: Variable Flow (Lens) ---
%         % 1. PyTorch -> PyBind
%         \draw[flowblue, boxstyle] (py_seq.north west) rectangle (py_seq.south east);
%         \draw[flowblue, boxstyle] (pb_seq.north west) rectangle (pb_seq.south east);
%         \draw[flowline, flowblue] (py_seq.north east) -- ++(1.5,0) |- (pb_seq.south west);

%         % 2. PyBind -> Host
%         \draw[flowblue, boxstyle] (host_lens.north west) rectangle (host_lens.south east);
%         \draw[flowline, flowblue] (pb_seq.south east) -- ++(1.8,0) |- (host_lens.north west);

%         % 3. Host -> Kernel (Argument)
%         \draw[flowblue, boxstyle] (kern_lens.north west) rectangle (kern_lens.south east);
%         \draw[flowline, flowblue] (host_lens.north east) -- ++(2,0) |- (kern_lens.south west);
        
%         % 4. Kernel Argument -> Usage
%         \draw[flowblue, boxstyle] (kern_lens_use.north west) rectangle (kern_lens_use.south east);
%         \draw[flowred, boxstyle] (kern_stage_spec.north west) rectangle (kern_stage_spec.south east);
%         % \draw[->, flowblue, dotted, thick] (kern_lens.south) -- ++(0,-0.3) -| (kern_lens_use.north);

%         % --- GREY ARROW: Type Mapping Relationship ---
%         % \draw[typearrow] (p1.south) -- node[left, font=\tiny, text=gray] {PyObject} (p_typemap_header.north -| p1.south);
%         % \draw[typearrow] (p_typemap_header.north -| p2.south) -- node[right, font=\tiny, text=gray] {C++ Struct} (p2.south);
%         \draw[typearrow] (p3.south -| p_llvm_header.south) -- node[right, font=\tiny, text=gray] {IR Generation} (p_llvm_header.north);
%         \draw[flowline, flowred] (p4.south -| p_spec_header.south) -- node[right, font=\tiny, text=flowred] {Specialization} (p_spec_header.north);

%     \end{tikzpicture}
%     \caption{Flash Attention Optimization Flow: Constant Propagation \& Field Reconstruction.}
% \end{figure}

